% !TEX root = ../main.tex

\section{Convergence Analysis of FedAMP}\label{sec:analysis}
% In this section, we analyze the convergence of FedAMP for both the situations when the function $\mG(W)$ in Eq.~\eqref{eq:our-form} is convex and non-convex, respectively. 
In this section, we analyze the convergence of FedAMP when $\mG$ is convex or non-convex under suitable conditions. 
%~\cite{rockafellar2009variational},
% The analysis of {FedAMP} is inspired by those in the optimization literature ({\it e.g.}, \cite{bertsekas2011incremental,li2019incremental}).
%For ease of exposition, we write $W^k = [w_1^k,\dots,w_m^k]$ and $U^k = [u_1^k,\dots,u_m^k]$ for every $k$. 
To begin with, similar to the analysis of many incremental and stochastic optimization algorithms~\cite{bertsekas2011incremental,nemirovski2009robust}, we make the following assumption.
\begin{assumption}\label{ass:bound}
There exists a constant $B>0$ such that $\max\{ \|Y\|: Y\in \partial \mF(W^k) \} \leq B$ and $\|\nabla \mA(W^k)\|\leq B/\lambda$ hold for every  $k \geq 0$, where $\partial \mF$ is the subdifferential of $\mF$ and $\|\cdot\|$ is the Frobenius norm. 
\nop{
\begin{equation}\nonumber
%\label{eq:ass-bound-1}
\max\{ \|Y\|: Y\in \partial \mF(w_1^k,\dots,w_m^k) \} \leq B
\end{equation}
and
\begin{equation}\nonumber
%\label{eq:ass-bound-2}
\|\nabla \mD(w_1^k,\dots,w_m^k)\|\leq B/\lambda.
\end{equation}
}
\end{assumption}

For our problem in Eq.~\eqref{eq:our-form}, Assumption~\ref{ass:bound} naturally holds if both $\mF(W)$ and $\mA(W)$ are locally Lipschitz continuous and $\|W^k\|$ is bounded by a constant for all $k\geq 0$. 

\nop{
\chu{bounded by what? a large enough finite constant?}.
}

Now, we provide the guarantee on convergence for FedAMP when both $\mF(W)$ and $\mA(W)$ are convex functions.
\nop{
When both $\mF$ and $\mD$ are convex, we have the following guarantee of {FedAMP}.
}
% \footnote{Note that for the case where $\mA$ is the tamed square root function defined in Section \ref{sec:instances}, the attention-inducing function $\mD$ is convex.}
\begin{thm}
\label{thm:analysis-convex}
Under Assumption \ref{ass:bound} and assuming functions $\mF(W)$ and $\mA(W)$ in Eq.~\eqref{eq:our-form} are convex, if $\alpha_{1} = \cdots = \alpha_K = \lambda/\sqrt{K}$ for some $K\geq 0$, then the sequence $W^0, \ldots, W^K$ generated by Algorithm \ref{alg:FedAMP} satisfies
\begin{equation}\nonumber
%\label{eq:analysis-convex}
\min_{0\leq k\leq K} \mG(W^k) \leq \mG^* + \frac{\|W^0 - W^*\|^2 + 5B^2}{\sqrt{K}},
\end{equation}
where $W^*$ is an optimal solution of Eq.~\eqref{eq:our-form} and $\mG^*= \mG(W^*)$. Moreover, if $\alpha_k$ satisfies $\sum_{k=1}^\infty\alpha_k = \infty$ and $\sum_{k=1}^\infty\alpha_k^2<\infty$, then %it holds that
\begin{equation}
\nonumber
\liminf_{k\rightarrow\infty} \mG(W^k) = \mG^*.
\end{equation}
\end{thm}

Theorem~\ref{thm:analysis-convex} implies that for any $\epsilon>0$, FedAMP needs at most $\mathcal{O}(\epsilon^{-2})$ iterations to find an $\epsilon$-optimal solution $\widetilde{W}$ of Eq.~\eqref{eq:our-form} such that $\mG(\widetilde{W}) - \mG^* \leq \epsilon$. It also establishes the global convergence of FedAMP to an optimal solution of Eq.~\eqref{eq:our-form} when $\mG$ is convex. The proof of Theorem \ref{thm:analysis-convex} is provided in Appendix~A~\cite{huang2020personalized}. 

\nop{\chu{To Zirui: Please write the remark of Theorem~\ref{thm:analysis-convex} in here, and also build the logical connection between Theorem~\ref{thm:analysis-convex} and our previous claim that ``$\mG(W)$ converges to an optimal point when it is convex''. Thanks!}}

\nop{
Recall that $\partial \mF(W) = \{\nabla\mF(W)\}$ if $\mF$ is differentiable at $W$.}

Next, we provide the convergence guarantee of {FedAMP} when $\mG(W)$ is a smooth and non-convex function. 

\begin{thm}
\label{thm:analysis-ncvx}
Under Assumption~\ref{ass:bound} and assuming functions $\mF(W)$ and $\mA(W)$ in Eq.~\eqref{eq:our-form} are continuously differentiable and the gradients $\nabla\mF(W)$ and $\nabla \mA(W)$ are Lipschitz continuous with modulus $L$, 
if $\alpha_{1} = \cdots = \alpha_K = \lambda/\sqrt{K}$, then the sequence $W^0, \ldots, W^K$ generated by Algorithm \ref{alg:FedAMP} satisfies
\begin{equation}\nonumber
%\label{eq:analysis-ncvx}
\begin{aligned}
& \min_{0\leq k\leq K} \|\nabla \mG(W^k)\|^2 \\ 
& \quad \leq \frac{18(\mG(W^0) - \mG^* + 20LB^2)}{\sqrt{K}} + \mathcal{O}\left(\frac{1}{K}\right)
% \mathcal{O}\left( \frac{\mG(w_1^0,\dots,w_m^0) - \mG^* + LB^2}{\sqrt{K}}\right) 
% \frac{34L\left(\mG(w^0,\dots,w^m) - \mG^* + 10L^2B^2\right)}{\sqrt{K}},
\end{aligned}
\end{equation}
where $W^*$ and $\mG^*$ are the same as in Theorem \ref{thm:analysis-convex}. Moreover, if $\alpha_k$ satisfies $\sum_{k=1}^\infty\alpha_k = \infty$ and $\sum_{k=1}^\infty\alpha_k^2<\infty$, then
\begin{equation}
\nonumber
\liminf_{k\rightarrow\infty} \|\nabla\mG(W^k)\| = 0.
\end{equation} 
\end{thm}

Theorem \ref{thm:analysis-ncvx} implies that for any $\epsilon>0$, FedAMP needs at most $\mathcal{O}(\epsilon^{-4})$ iterations to find an $\epsilon$-approximate stationary point $\widetilde{W}$ of Eq.~\eqref{eq:our-form} such that $\|\nabla\mG(\widetilde{W})\|\leq \epsilon$. It also establishes the global convergence of FedAMP to a stationary point of Eq.~\eqref{eq:our-form} when $\mG$ is smooth and non-convex. The proof of Theorem~\ref{thm:analysis-ncvx} is in Appendix~B~\cite{huang2020personalized}.

\nop{\chu{To Zirui: Please write the remark of Theorem~\ref{thm:analysis-ncvx} in here, and also build the logical connection between Theorem~\ref{thm:analysis-ncvx} and our previous claim that ``$\mG(W)$ converges to a stationary point when it is non-convex''. Thanks!}}
